% ============================================================================
% Fichier  : partie2/P02_C05_normes_referentiels.tex
% Titre    : Normes et Référentiels de l'Auditeur Public
% Auteur   : MINKA MI NGUIDJOÏ Thierry Emmanuel
% Version  : 0.1 -- novembre 2026
% ============================================================================
% Sources : INF3226 C13 et C14 (resserrés) ; GUID 5100 et manuel IDI
% (dossier Normes) ; normes IIA 2017 (DOC-20230213, DOC-20330213) ; COBIT
% 2019 ; ISO/IEC 27001 et 27002 (2022) ; lois n° 2010/012 et 2024/017.
% Aucune norme ISO ou COBIT n'est reproduite (droits) : résumés seulement.
% ============================================================================

\chapter{Normes et Référentiels de l'Auditeur Public}
\label{chap:normes}
\chaptermark{Normes et référentiels}

\epigraph{%
    « La norme dit comment auditer~; le référentiel dit à quoi comparer. »%
}{Règle du cours}

\niveauChapitre{Fondamental}{Chapitre~\ref{chap:institutions}.}

\begin{boxobjectifs}
\begin{itemize}
    \item Distinguer une norme professionnelle (comment auditer) d'un
          référentiel de critères (à quoi comparer).
    \item Situer les normes de l'INTOSAI (ISSAI, GUID~5100) et le manuel de
          l'IDI dans le travail du vérificateur.
    \item Situer les normes de l'IIA dans le travail de l'auditeur interne.
    \item Utiliser COBIT~2019 et ISO/IEC~27001 et~27002 comme sources de
          critères, sans les confondre avec la loi.
    \item Choisir, pour un test donné, un critère explicite et le citer
          correctement.
    \item Énoncer les principes de déontologie communs aux auditeurs.
\end{itemize}
\end{boxobjectifs}

\bigskip

% ============================================================================
\section*{Pourquoi ce chapitre}
% ============================================================================

Le prologue l'a montré~: une cote \NON{} n'a de sens que si l'on sait à
quoi l'on compare. Ce chapitre sépare deux choses que l'on confond
souvent. Les \emph{normes professionnelles} disent comment conduire la
mission~; les \emph{référentiels} et les \emph{textes} fournissent les
critères. L'auditeur public a besoin des deux.

\separateur

% ============================================================================
\section{Norme, référentiel, critère}
\label{sec:c05-notions}
% ============================================================================

\begin{boxdef}
\termecle{Norme professionnelle}\index{norme professionnelle}~: règle qui
encadre la conduite d'une mission~: indépendance, planification, preuve,
documentation, rapport. Elle s'impose à l'auditeur.

\medskip
\termecle{Référentiel}\index{référentiel}~: ensemble structuré de bonnes
pratiques ou d'exigences (COBIT, ISO/IEC~27002) dans lequel l'auditeur
puise des critères. Il ne s'impose à l'entité que si un texte ou un
engagement le rend obligatoire.

\medskip
\termecle{Critère d'audit}\index{critère d'audit}~: exigence de référence,
choisie avant le test, à laquelle on compare la situation constatée.
\end{boxdef}

\begin{tableaucours}{L{3.0cm}L{4.4cm}L{4.8cm}}{%
    \tableauHeader{Source} &
    \tableauHeader{Force} &
    \tableauHeader{Exemple de critère}
}
Loi, décret & Obligatoire & Les SI des organismes publics font l'objet
d'un audit de sécurité (loi n°~2010/012, art.~32) \\
Règle interne & Obligatoire pour l'entité & Procédure de la DAF sur la
désactivation des comptes \\
Référentiel adopté & Obligatoire si l'entité l'a adopté & Politique de
sécurité fondée sur ISO/IEC~27002 \\
Bonne pratique & Recommandation & Un objectif de gestion COBIT~2019 \\
\end{tableaucours}
\vspace{-0.8em}
\begin{center}{\small\itshape Tableau~5.1. Hiérarchie des sources de
critères}\end{center}

\begin{boxalert}
\textbf{Violation typique.} Écrire «~non conforme à ISO/IEC~27002~» pour
une administration qui n'a jamais adopté cette norme.

\smallskip
\textbf{Règle~:} on n'est «~non conforme~» qu'à une règle qui s'impose.
Face à une simple bonne pratique, on écrit «~l'entité ne dispose pas de…,
alors que les bonnes pratiques (référence) recommandent…~», et l'on
classe l'écart en conséquence.
\end{boxalert}

\separateur

% ============================================================================
\section{Les normes de l'INTOSAI}
\label{sec:c05-intosai}
% ============================================================================

\subsection{Le cadre}

L'INTOSAI réunit les institutions supérieures de contrôle des finances
publiques, dont le CONSUPE fait partie au titre de l'article~2 du décret
n°~2013/287. Son cadre de prises de position professionnelles comprend
des principes, des normes (\sigle{ISSAI}) et des lignes directrices
(\sigle{GUID}).

\begin{tableaucours}{L{2.6cm}L{9.6cm}}{%
    \tableauHeader{Texte} &
    \tableauHeader{Objet}
}
ISSAI~100 & Principes fondamentaux de l'audit du secteur public~: les
trois types d'audit, les éléments de l'audit, les principes généraux \\
ISSAI~200, 300, 400 & Principes de l'audit financier, de performance et de
conformité \\
ISSAI~4000 & Norme de l'audit de conformité \\
GUID~5100 & Lignes directrices sur la vérification des systèmes
d'information \\
Manuel de l'IDI & Mise en œuvre pratique de l'audit de conformité selon les
ISSAI \\
\end{tableaucours}
\vspace{-0.8em}
\begin{center}{\small\itshape Tableau~5.2. Les textes de l'INTOSAI utiles à
l'auditeur des SI publics}\end{center}

\subsection{Le GUID 5100}

Le GUID~5100 complète les principes de l'ISSAI~100 pour les missions qui
portent sur les systèmes d'information~\cite{intosai2019guid5100}. Il
précise qu'un audit des SI peut être autonome ou s'insérer dans un audit
financier, de conformité ou de performance, et il énumère treize domaines
que la mission peut couvrir, parmi lesquels la politique et la
gouvernance des TI, le développement et l'acquisition des systèmes, la
gestion des opérations, la sécurité de l'information, la conformité
réglementaire, la continuité et la reprise après sinistre, et les
contrôles d'applications. Le cours suit la même logique~:
la Partie~IV reprend ces domaines.

\subsection{L'audit de conformité selon l'IDI}

Le manuel de l'IDI décrit l'audit de conformité comme une relation entre
trois parties~: l'auditeur, la partie responsable (l'entité auditée) et
les utilisateurs du rapport~\cite{idi2021conformite}. Il insiste sur deux
exigences que l'auditeur des SI applique directement~: identifier des
\emph{critères} appropriés avant le travail de terrain, et documenter
suffisamment la mission pour qu'un auditeur expérimenté, sans lien avec
elle, comprenne les travaux et les conclusions.

\separateur

% ============================================================================
\section{Les normes de l'audit interne}
\label{sec:c05-iia}
% ============================================================================

L'auditeur interne d'un ministère ou d'un établissement public suit les
normes de l'Institut des auditeurs internes (\sigle{IIA}). Le dossier des
normes contient l'édition 2017 du Cadre de référence international des
pratiques professionnelles~\cite{iia2017cripp}~; l'IIA l'a remplacée par
les normes mondiales d'audit interne publiées en
2024~\cite{iia2024gias}. Les exigences qui comptent ici sont stables
d'une édition à l'autre~: indépendance et objectivité, compétence,
planification fondée sur les risques, preuve suffisante, communication
des résultats et suivi des recommandations.

\begin{boxinfo}
\textbf{Vérificateur et auditeur interne parlent la même langue.} Les
normes de l'INTOSAI et celles de l'IIA convergent sur l'essentiel~:
critère, preuve, documentation, contradictoire. C'est ce qui permet au
CONSUPE de s'appuyer sur les rapports d'audit interne qu'il reçoit au
titre de l'article~5 du décret n°~2013/287.
\end{boxinfo}

\separateur

% ============================================================================
\section{Les référentiels de critères}
\label{sec:c05-referentiels}
% ============================================================================

\subsection{COBIT 2019}

COBIT~2019, de l'ISACA, est un référentiel de gouvernance et de gestion
des informations et des technologies~\cite{isaca2018cobitintro}. Il
organise quarante objectifs en cinq domaines~: un domaine de gouvernance
(évaluer, diriger, surveiller) et quatre domaines de
gestion~\cite{isaca2018cobitobjectifs}. L'auditeur public y puise des
critères pour la gouvernance de la fonction informatique, les projets, et
la gestion des changements (Chapitres~\ref{chap:exploitation}
et~\ref{chap:gouvernance}).

\subsection{ISO/IEC 27001 et 27002}

ISO/IEC~27001 fixe les exigences d'un système de management de la
sécurité de l'information~; ISO/IEC~27002 détaille les mesures de
sécurité, regroupées dans l'édition 2022 en quatre thèmes~: mesures
organisationnelles, liées aux personnes, physiques et
technologiques~\cite{iso27001_2022,iso27002_2022}. Ces normes sont sous
droits~: le cours les résume, il ne les reproduit pas.

\subsection{Les textes camerounais comme critères}

\begin{tableaucours}{L{3.4cm}L{8.8cm}}{%
    \tableauHeader{Texte} &
    \tableauHeader{Ce qu'il fournit comme critère}
}
Décret n°~2013/287 & Le droit du vérificateur à obtenir les pièces,
électroniques comprises \\
Loi n°~2010/012 & L'audit de sécurité des SI des organismes publics, les
obligations de sécurité des exploitants \\
Loi n°~2018/012 (régime financier) & Les principes de l'exécution de la
dépense et la séparation ordonnateur / comptable (à lire dans le
texte)~\cite{minfi2023regimefinancier} \\
Loi n°~2024/017 & La protection des données personnelles traitées par
les SI publics (paie, contribuables, usagers) \\
\end{tableaucours}
\vspace{-0.8em}
\begin{center}{\small\itshape Tableau~5.3. Textes camerounais utiles comme
critères}\end{center}

\subsection{Choisir et citer un critère}

\begin{boxprocedure}
\textbf{Fiche de critère}, une par test~:
\begin{itemize}
    \item l'énoncé du critère, en une phrase vérifiable~;
    \item sa source exacte (texte, article, alinéa~; ou référentiel et
          objectif), et sa force (Tableau~5.1)~;
    \item la date à laquelle le texte a été lu, et dans quelle version~;
    \item le test qui permettra de dire \OUI{} ou \NON{}.
\end{itemize}
\end{boxprocedure}

\begin{boxscenario}{Dossier MINÉQUIR. Le critère du contrôle 1}
\textbf{Constatation.} La DAF n'a aucune procédure écrite sur la
désactivation des comptes.

\medskip
\textbf{Analyse} (confiance élevée). Faute de règle interne, la
vérificatrice fonde son critère sur une bonne pratique
(principe~11 du COSO~; gestion des accès dans ISO/IEC~27002) et le
signale comme tel. Son constat portera donc sur deux points~: l'absence
de règle interne, qui est elle-même une déficience de conception, et les
9 comptes encore actifs, qui en sont la conséquence.
\end{boxscenario}

\separateur

% ============================================================================
\section{Déontologie}
\label{sec:c05-deontologie}
% ============================================================================

Les codes de l'INTOSAI, de l'IIA et de l'ISACA convergent sur quatre
exigences.

\begin{tableaucours}{L{3.0cm}L{9.2cm}}{%
    \tableauHeader{Exigence} &
    \tableauHeader{Ce qu'elle impose dans une mission SI}
}
Intégrité & Ne jamais modifier une donnée de l'entité~; travailler sur
une copie \\
Indépendance et objectivité & Ne pas auditer un système que l'on a
conçu, paramétré ou exploité \\
Compétence & Se faire assister par un spécialiste pour les tests
techniques que l'on ne maîtrise pas, et le dire au rapport \\
Confidentialité & Protéger les extractions (paie, contribuables) et les
détruire selon les règles à la fin de la mission \\
\end{tableaucours}
\vspace{-0.8em}
\begin{center}{\small\itshape Tableau~5.4. Quatre exigences
déontologiques dans une mission SI}\end{center}

\begin{boxmetier}{Vérificateur}
Écrivez dans le programme de travail, pour chaque test, la fiche de
critère. Lors du contradictoire, la première contestation porte presque
toujours sur le critère~: «~aucun texte ne nous l'impose~».
\end{boxmetier}

\begin{boxmetier}{Auditeur interne}
Si votre ministère a adopté une politique de sécurité fondée sur
ISO/IEC~27002, ses mesures deviennent des règles internes~: citez la
politique, pas la norme.
\end{boxmetier}

\separateur

\begin{boxresume}
\begin{itemize}
    \item \textbf{Norme et référentiel.} La norme dit comment auditer~; le
          référentiel fournit des critères.
    \item \textbf{L'INTOSAI pour le vérificateur.} ISSAI~100, 400 et 4000,
          GUID~5100, manuel de l'IDI.
    \item \textbf{L'IIA pour l'auditeur interne.} Mêmes exigences
          essentielles~: critère, preuve, documentation.
    \item \textbf{COBIT et ISO/IEC~27002 comme sources de critères}, qui ne
          s'imposent que si l'entité les a adoptés.
    \item \textbf{Une hiérarchie des critères.} Loi, règle interne,
          référentiel adopté, bonne pratique~; l'écart se qualifie selon la
          force du critère.
    \item \textbf{Une fiche de critère par test}, avec la source exacte et la
          date de lecture.
    \item \textbf{Quatre exigences déontologiques.} Intégrité,
          indépendance, compétence, confidentialité.
\end{itemize}
\end{boxresume}

\begin{boxexo}[title={\ding{45}\quad Exercice~5.1. Norme ou référentiel~? \niveaubase}]
Classez~: (a) GUID~5100~; (b) ISO/IEC~27002~; (c) ISSAI~4000~; (d) COBIT~2019~;
(e) normes de l'IIA~; (f) loi n°~2024/017.
\end{boxexo}

\begin{boxexo}[title={\ding{45}\quad Exercice~5.2. Qualifier l'écart \niveaumoyen}]
La CAPNO n'a pas testé ses sauvegardes depuis deux ans. Rédigez le
constat dans deux cas~: (a) la politique de sécurité de la CAPNO exige un
test semestriel~; (b) aucune règle interne n'existe.
\end{boxexo}

\begin{boxexo}[title={\ding{45}\quad Exercice~5.3. Fiches de critère \niveaumoyen}]
Établissez la fiche de critère des contrôles~2 (séparation LIQ / ORD) et~4
(blocage des doublons) du circuit de la dépense du MINÉQUIR
(Tableau~3.2).
\end{boxexo}

\begin{boxexo}[title={\ding{45}\quad Exercice~5.4. Un conflit d'intérêts \niveauexpert}]
Un vérificateur a participé, deux ans plus tôt, au comité de pilotage du
projet SIDÉPENSE. Analysez la situation au regard des exigences
déontologiques, et proposez ce que son chef de division doit décider.
\end{boxexo}

\begin{boxinfo}
\textbf{Transition.} Le mandat et les normes sont posés. La Partie~III
conduit la mission elle-même, en commençant par sa préparation
(Chapitre~\ref{chap:preparer}).
\end{boxinfo}
